\section{The global index with all ramified places declared as rows}

% BEGIN I81G1
\paragraph{I81G1. A unit root difference on the ramified hub.}
Use the member, good localization, point and complement of H80M1.
Write $p=\ell\in\{q,d\}$, $K=E_w$, and $K_0$ for the maximal
unramified subfield of $K$. At every $w\mid p$, R79Q2 gives
$K\subset K_0(\mu_p)$ with $[K:K_0]=2$. Since
$K_0(\mu_p)/K_0$ is totally ramified and cyclic of order $p-1$,
$K(\mu_p)/K$ is totally ramified and cyclic of order $(p-1)/2$.
Lift a generator to local inertia, then to the specified global
inertia group through the place identification. For its element
$g_0$, the residues of $\chi_c(g_0)$, $1\le c\le t$, are distinct.
Indeed they equal $\chi_1(g_0)\varepsilon_p(g_0)^{-(c-1)}$
and $t<(p-1)/2$ by the host bound $p>4tr+1$.

D79G1--D79G3 apply with their original hypotheses at this localization.
The polynomial of the character corner at $g_0$ is monic of degree
$t$ over the local Hecke algebra. At the chosen point it has the
$t$ simple residue roots just found. The Hecke coefficients have
the same residues at all points. The Hensel argument in D79G4
therefore supplies the $t$ separately labelled roots at all points.
D79G5 and the construction in D79G6 now give the ordered-root cover,
the integral projectors and the unit $\delta(\pi')$ used in G79E7. The general decomposition
record also asks for residue order $p-1$. Here the order is
$(p-1)/2$. The proofs of G79E7 and G79E11--G79E13 below use
only the listed integral blocks, root basis and unit difference.
The off-support and away-from-$p$ hypotheses are supplied by D79G7.
Thus G79E11--G79E12 give, with the actual point ideal and block
cocycle modules of G79E1 and G79E7,
\[
  \operatorname{length}_O(O/I_\pi)\le
  \operatorname{length}_O C_\rho+
  \sum_{0\le j<t}\operatorname{length}_O C_j .
\]
Each intersects the global coboundaries in zero. Moreover the column
construction of G79E13 puts the restriction of every generator in
the local cut set for the corresponding block. That construction
requires the integral projector and the intertwining map on its
columns. Both have just been supplied at $q,d$.
% END I81G1

% BEGIN I81G2
\paragraph{I81G2. Normalized lengths after enlarging coefficients.}
Use the specified bridge to identify the coefficient DVR of a
localization with $O$ at $w$. Choose any good finite coefficient
extension with rings $O\subset O'$ and fields $L\subset L'$.
The given bridge field already contains all roots of unity of order
$\operatorname{lcm}(p,\mathrm{condN})$ and splits $K$. These two
properties persist under every finite coefficient enlargement. The test over
$G_K$ keeps $K$ fixed. For a lattice $T\subset V$ put $A=V/T$.
Writing $O'$ in a free $O$-basis identifies the continuous cochains
of $T'=T\otimes_O O'$, $V'=V\otimes_L L'$ and $A'=V'/T'$
with finite direct sums of the original cochains. This identification
commutes with the differential, and so with cohomology. For $V$,
being Bloch--Kato finite is the existence of a cobounding primitive
after tensoring with the specified rational crystalline period ring.
In the basis above this is exactly the corresponding assertion in
each coordinate. For $T$ take the inverse image of the rational finite
part, and for $A$ take its image. We obtain the same base change
for finite parts and for $H^1(K,A)/H_f^1(K,A)$.

Consequently the $O'$-span in this last quotient after enlargement
of any given set of original classes is its original $O$-span
tensored with $O'$. This follows by flatness, including for any
finite subset. If $d=e(O'/O)$, a finite-length module of length
$h$ has extended length $dh$, by extending a composition series
and using $\varpi O'=(\varpi')^d O'$. The same factor multiplies
$\operatorname{length}_O(O/p)$. A uniform finite bound on the
normalized lengths over $O'$ therefore applies over $O$. To see
this without any finite-generation assumption on the full span,
apply it first to finite subsets. A strictly increasing chain of
finitely generated submodules cannot have more steps than the
uniform length bound. Starting with zero and adjoining a generator
outside the current span must therefore exhaust the span.
% END I81G2

% BEGIN I81G3
\paragraph{I81G3. One integral row map for any finite list of rank cuts.}
Work over $O'$ and write $V$ for the enlarged rational rank block
$\rho/\chi_1$, with its original extended lattice $T$. By C80L1
each rank cut $z_j$ in a given finite list has the presentation
$z_j(g)=\phi_j((g-1)E_j)$ on a stable finite submodule $Y_j$
in a sum $P_j$ of actual $V_x=R_x/\chi_1$. Take
$P=\bigoplus P_j$, $Y=\bigoplus Y_j$ and
$\phi=\sum\phi_j:Y\to V/T$. Choose a nonzero nonunit $c\in O'$
killing $\phi(Y)$, possible by finite generation. Choose a full
stable lattice $\Lambda$ of $P$, using the original integral
matrices in each summand. The submodules of $\Lambda$ defined by
$c^h v\in Y$ form an ascending chain and hence stabilize, say at
$h=n$. Then
$Y\cap c^{n+1}\Lambda\subset cY$: write a vector as $c^{n+1}v$
and apply stabilization to $v$. On $L_Y=Y+c^{n+1}\Lambda$
extend $\phi$ by zero on $c^{n+1}\Lambda$. This gives a
well-defined equivariant map from a full lattice, also when
$Y$ has smaller rational rank.

In $P\oplus V$ use the stable lattice
\[
 U=\{(y,v): y\in L_Y,\quad v\bmod T=\phi(y)\}.
\]
Projection onto $L_Y$ has kernel $0\oplus T$ and splits over $O'$
by lifting a basis. Thus $U$ is finite free and spans $P\oplus V$,
and the inclusion of $T$ induces an injection
$j:V/T\longrightarrow (P\oplus V)/U$.
For $d_{j,g}=(g-1)E_j$ in the $j$th summand of $Y$,
$(d_{j,g},\phi(d_{j,g}))$ interpreted with any lift in $V$
lies in $U$. It follows that
\[
 j(z_j(g))=(-d_{j,g},0)\bmod U
           =(g-1)(-E_j,0)\bmod U .
\]
Every class in the list is therefore killed by the same map on
local cohomology. Restricting integral functionals along $T\to U$
gives a surjective map $F:U^\vee(1)\to T^\vee(1)$
by the free splitting above. The rational source embeds
equivariantly in a sum of the permitted $V_x^\vee(1)$:
$P$ is already such a sum before dualizing, and $V$ is a direct
summand of $V_\pi$ by the specified good-point splitting.
The target with its lattice is exactly the matrix-lattice dual
in the row declaration. The representations and maps are continuous
in their finite-dimensional topologies. Hence every one of these
maps lies in the quantified class defining $\mathrm{RowAt}$.
% END I81G3

% BEGIN I81G4
\paragraph{I81G4. The rank defect is bounded by this single row length.}
Suppose the maps of I81G3 have normalized row length at most
$B_w$ for all finite lists at $w$. The full row bound is sufficient.
The quotient measured by $F.\mathrm{col}$ is
\[
 D_F=H_f^1(K,T^\vee(1))/F_* H_f^1(K,U^\vee(1)).
\]
Indeed the cocycle definition quotients the finite integral cocycles
by images of finite cocycles and target integral coboundaries.
The Tate pairing of $H^1(K,V/T)$ with $H^1(K,T^\vee(1))$
annihilates $F_* H_f^1(K,U^\vee(1))$ on all the classes in the
list. Pull the pairing back along $T\to U$ by functoriality:
the first argument becomes a coboundary by I81G3. The annihilator
of the full $H_f^1(K,T^\vee(1))$ is precisely $H_f^1(K,V/T)$
by Bloch--Kato local orthogonality for the stated de Rham pair.
This is the torsion--integral finite-part comparison on the tied
period rings in the local Bloch--Kato and Tate hypotheses. The
integral dual can equivalently be computed after restriction to
$\mathbb Z_p$. Choose a generator $b$ of the inverse different
of $O'$ over $\mathbb Z_p$ and use $x\mapsto\operatorname{Tr}(bx)$.
This identifies $O'$ with its full integral trace dual and induces
the perfect pairing with the discrete quotient. Since the tested
finite subgroup is an $O'$-module, testing after this trace detects
every scalar pairing (test all its $O'$-multiples).

The span of the classes modulo the finite part thus injects in
the Pontryagin dual of $D_F$. For a finite $O'$-module duality
preserves length, by a composition series and duality on the
finite residue field. Hence this span has length at most
$B_w\operatorname{length}_{O'}(O'/p)$, regardless of the number
of generators or their individual presentations. I81G2 gives
the same normalized bound on the entire original rank cut span.
% END I81G4

% BEGIN I81G5
\paragraph{I81G5. The length lost on each character block.}
Write $A_e=(L'/O')(-e)$ for the quotient of the integral cyclotomic
line of exponent $-e$, $0\le e<t$. For one character cut take
$P,E,Y,\phi$ as in C80L1. In $P\oplus L'$ with trivial last
summand take the lattice $U=Y\oplus O'(E,1)$ and its rational
span. The sum is direct by the last coordinate, and it is stable
since $(g-1)E\in Y$. Its rational representation embeds into a sum
of the $V_x$, using the trivial direct summand of $V_\pi$. This
representation is crystalline. Its labelled jumps are contained
in the common jump sets of $V_\pi$ by strictness and S80J1.
Projection to the last coordinate is invariant and integral on $U$,
takes $(E,1)$ to $1$ and has integral kernel exactly $Y$.
Extend $\phi$ to $U$ by zero on $(E,1)$ and then $O'$-linearly
to its rational span: $A_e$ is divisible and injective over the
DVR, since a map on a nonzero principal ideal extends by dividing
its generator value. The equivariance tested by $\mathrm{TBAt}$
is only on the integral kernel, where it holds by construction.

Consequently the first clause of $\mathrm{TBAt}$ at $c_B$ puts
$p^{c_B} z$ in the finite part for every original cut. By linearity
this holds on their span at the enlarged coefficients.
H80M4 bounds its length modulo the finite part by
$(d_w+1)c_B\operatorname{length}_{O'}(O'/p)$,
where $d_w=[K:\mathbb Q_p]$. By I81G2 the same normalized length
bound holds over the original coefficients. Above $p$, finite
classes in the trivial ratio are unramified and those in the
nontrivial ratios are zero as cohomology classes, as in H80M4.
% END I81G5

% BEGIN I81G6
\paragraph{I81G6. The all-row inequality on the original Hecke ideal.}
Fix the all-row declaration of R79Q4. Suppose at each $w\mid p$
that $\mathrm{TBAt}$ holds at $c_B$ and that the maps of I81G3
satisfy the row length bound $B_w$. Write $e_O=\operatorname{length}_O(O/p)$.
Restrict $C_\rho$ simultaneously to the local quotients by $H_f^1$
at $w\mid p$. By I81G1 and I81G4 its image has length at most
$e_O\sum_w B_w$. Its kernel maps injectively to the native Selmer
group: at $p$ it is finite by definition of the kernel, away from
$S_i,Q_i,p$ inertia vanishes by D79G7 and G79E8, and G79E7
excludes global coboundaries. All these maps use the same global
lattice and point.

For the $e$th character, first restrict to inertia at $Q_i$, losing
at most $\sum_{P\in Q_i}z_P$ by G79E8, with $z_P$ as in G79E1.
The kernel vanishes on all inertia away from $p$ by D79G7. On
this kernel restrict to the local quotients by the finite parts
at $w\mid p$. Its image has length at most
$e_O c_B\sum_w(d_w+1)$ by I81G5. The new kernel for $e=0$
consists of everywhere unramified homomorphisms, hence has length
at most $e_O v_p(h_E)$ by Artin reciprocity as in G79E6.
For $0<e<t$ the new kernel is locally a coboundary above $p$,
hence has length at most $e_O u_e$ by G79E6--G79E7.
Under the original finiteness guards all these bounds are finite.
Add them in I81G1 and divide by $e_O$ using G79E1, to get
\[
\begin{split}
 v_{\rm CC}(i,\pi)\ \le\;&
 s_\pi+\sum_{0<e<t}u_e+v_p(h_E)
       +\frac{t}{e_O}\sum_{P\in Q_i}z_P\\
 &+\sum_{w\mid p}B_w+t c_B\sum_{w\mid p}(d_w+1).
\end{split}
\]
This implies the cited $\mathrm{GlobalLatticeIndexTheorem}$ at
these twisted indices at the two primes, for each good point $\pi$
and its given complement, with every place a row. In the native expression
the character coefficient is $d_w+2$, the empty column set
contributes zero to its column sums, and the additional term
$4n_Vv_\delta(i)$ is nonnegative by R79Q4. The killing inequality
bounds the auxiliary-prime correction by $p^{-2}$ on its original
guard $tr<[E(\mu_p):E]$. This proves the required implication
from the local estimates on this hub directly for its evaluation
ideal, with no global index inequality left as a premise.
% END I81G6
